
\documentclass{article}
\usepackage{colortbl}
\usepackage{makecell}
\usepackage{multirow}
\usepackage{supertabular}

\begin{document}

\newcounter{utterance}

\centering \large Interaction Transcript for game `cladder', experiment `full\_v1.5\_default', episode 7073 with llama{-}3.1{-}8b{-}instruct.
\vspace{24pt}

{ \footnotesize  \setcounter{utterance}{1}
\setlength{\tabcolsep}{0pt}
\begin{supertabular}{c@{$\;$}|p{.15\linewidth}@{}p{.15\linewidth}p{.15\linewidth}p{.15\linewidth}p{.15\linewidth}p{.15\linewidth}}
    \# & \multicolumn{2}{c}{Player} && \multicolumn{2}{c}{Game Master} \\
    \hline

    \theutterance \stepcounter{utterance}  
    & & & \multicolumn{4}{p{0.6\linewidth}}{
        \cellcolor[rgb]{0.9,0.9,0.9}{
            \makecell[{{p{\linewidth}}}]{
                \texttt{\tiny{[P1$\langle$GM]}}
                \texttt{You are an expert in causal inference. The following question is not a typical commonsense query, but rather a meticulously designed question created by a professor specializing in causal inference, intended to assess the students' mastery of the course content.} \\
\texttt{Imagine a self{-}contained, hypothetical world with only the following conditions, and without any unmentioned factors or causal relationships: Jyka has a direct effect on yupt. Yupt has a direct effect on kwox. The overall probability of jyka is 25\%. The probability of not jyka and kwox is 71\%. The probability of jyka and kwox is 24\%. Is the chance of kwox smaller when observing jyka?} \\
\texttt{Start your answer with 'Yes' or 'No', followed by additional reasoning or evidence to support your explanation.} \\
            }
        }
    }
     \\ \\

    \theutterance \stepcounter{utterance}  
    & \multicolumn{4}{p{0.6\linewidth}}{
        \cellcolor[rgb]{0.561,0.737,0.561}{
            \makecell[{{p{\linewidth}}}]{
                \texttt{\tiny{[P1$\rangle$GM]}}
                \texttt{No.} \\
\\ 
\texttt{To determine whether the chance of kwox is smaller when observing jyka, we need to calculate the conditional probability of kwox given jyka (P(kwox|jyka)) and compare it to the unconditional probability of kwox (P(kwox)).} \\
\\ 
\texttt{We are given the following information:} \\
\\ 
\texttt{1. P(jyka) = 0.25} \\
\texttt{2. P(not jyka and kwox) = 0.71} \\
\texttt{3. P(jyka and kwox) = 0.24} \\
\\ 
\texttt{We can use the law of total probability to calculate P(kwox):} \\
\\ 
\texttt{P(kwox) = P(kwox and jyka) + P(kwox and not jyka)} \\
\texttt{= P(jyka and kwox) + P(not jyka and kwox)} \\
\texttt{= 0.24 + 0.71} \\
\texttt{= 0.95} \\
\\ 
\texttt{Now, we need to calculate P(kwox|jyka):} \\
\\ 
\texttt{P(kwox|jyka) = P(kwox and jyka) / P(jyka)} \\
\texttt{= 0.24 / 0.25} \\
\texttt{= 0.96} \\
\\ 
\texttt{Comparing P(kwox|jyka) and P(kwox), we see that P(kwox|jyka) is actually larger than P(kwox). This means that observing jyka increases the chance of kwox, rather than decreasing it.} \\
            }
        }
    }
    & & \\ \\

    \theutterance \stepcounter{utterance}  
    & & & \multicolumn{2}{p{0.3\linewidth}}{
        \cellcolor[rgb]{0.9,0.9,0.9}{
            \makecell[{{p{\linewidth}}}]{
                \texttt{\tiny{[GM$|$GM]}}
                \texttt{No} \\
            }
        }
    }
    & & \\ \\

    \theutterance \stepcounter{utterance}  
    & & & \multicolumn{2}{p{0.3\linewidth}}{
        \cellcolor[rgb]{0.9,0.9,0.9}{
            \makecell[{{p{\linewidth}}}]{
                \texttt{\tiny{[GM$|$GM]}}
                \texttt{no} \\
            }
        }
    }
    & & \\ \\

    \theutterance \stepcounter{utterance}  
    & & & \multicolumn{2}{p{0.3\linewidth}}{
        \cellcolor[rgb]{0.9,0.9,0.9}{
            \makecell[{{p{\linewidth}}}]{
                \texttt{\tiny{[GM$|$GM]}}
                \texttt{game\_result = WIN} \\
            }
        }
    }
    & & \\ \\

\end{supertabular}
}

\end{document}
